\documentclass[12pt, a4paper]{article}
\usepackage[utf8]{inputenc}
\usepackage[T1]{fontenc}
\usepackage[english]{babel}
\usepackage{geometry}
\geometry{top=2.5cm, bottom=2.5cm, left=2cm, right=2cm}
\usepackage{parskip}
\usepackage{microtype}
\usepackage[default]{sourcesanspro}
\usepackage{xcolor}
\usepackage{fancyhdr}
\usepackage[most]{tcolorbox}
\usepackage[hidelinks]{hyperref}
\usepackage{amssymb, amsmath}

\newcommand{\Book}{Principles of Mathematical Analysis}
\newcommand{\Chapter}{The real and complex number systems}
\definecolor{StructureColor}{RGB}{40, 80, 120}

\pagestyle{fancy}
\fancyhf{}
\lhead{\color{StructureColor}\textbf{\Book}}
\rhead{\small\Chapter}
\cfoot{\thepage}
\renewcommand{\headrulewidth}{0.4pt}
\renewcommand{\headrule}{\hbox to\headwidth{\color{StructureColor}\leaders\hrule height \headrulewidth\hfill}}

\newtcolorbox{exercisebox}[1]{
    enhanced, colback=StructureColor!5!white, colframe=StructureColor,
    coltitle=white, fonttitle=\bfseries\sffamily,
    title={#1}, sharp corners, boxrule=0.5mm
}

\begin{document}
\thispagestyle{fancy}

\begin{exercisebox}{Exercise 2}
    Prove that there is no rational number whose square is 12.
\end{exercisebox}

\textbf{Solution}

Suppose, for contradiction, that there exists a rational number $x$ such that $x^2=12$.

Note that
\[
\begin{aligned}
    x^2&=12\\
    x&=\sqrt{12}\\
    x&=\sqrt{4\cdot 3}\\
    x&= 2\sqrt{3}
\end{aligned}
\]
Since $\sqrt{3}$ is irrational (Proposition), which implies that $x$ is also irrational, this is a contradiction .

Therefore, there is no rational number whose square is 12. $\blacksquare$

\newpage

\begin{exercisebox}{Proposition}
    $\sqrt{3}$ is irrational.
\end{exercisebox}

\textbf{Solution}

Suppose, for contradiction, that $\sqrt{3}$ is a rational number. Then exists $a,b\in \mathbb{Z}$ with $a,b\neq 0$ such that
\[
\sqrt{3} = \dfrac{a}{b},
\]
where $a$ and $b$ are coprime. So,
\[
\begin{aligned}
    \sqrt{3}\cdot b &= a
\end{aligned}
\]
if b is even.

Then $b=2\cdot k$. Note that
\[
\begin{aligned}
    \sqrt{3}\cdot b &= a\\
    \sqrt{3} \cdot 2 \cdot k &= a\\
    2 \cdot (k \cdot \sqrt{3}) &= a
\end{aligned}
\]

Note that $a$ and $b$ are even, this is a contradiction.

if b is odd.

Then $b=2\cdot k + 1$. Note that

\[
\begin{aligned}
    \sqrt{3}\cdot b &= a\\
    \sqrt{3} \cdot (2\cdot k + 1) &= a\\
     (\sqrt{3} \cdot (2\cdot k + 1))^2&= a^2\\
     3 \cdot (4\cdot k^2+4 \cdot k+1)&=a^2\\
     12 \cdot k^2 + 12 \cdot k + 3&= a^2\\
     12 \cdot k^2 + 12 \cdot k + 2 + 1 &= a^2\\
     2(6 \cdot k^2 + 6 \cdot k + 1) + 1 &= a^2
\end{aligned}
\]

if a is even

Then $a=2\cdot q$. Note that
\[
\begin{aligned}
     2\cdot(6 \cdot k^2 + 6 \cdot k + 1) + 1&= a^2\\
     2\cdot(6 \cdot k^2 + 6 \cdot k + 1) + 1&= (2\cdot q)^2\\
     2\cdot(6 \cdot k^2 + 6 \cdot k + 1) + 1&= 4\cdot q^2\\
     2\cdot(6 \cdot k^2 + 6 \cdot k + 1) + 1 &= 2 \cdot (2\cdot q^2)\\
\end{aligned}
\]

The left side is a odd number and the right side is even number, this is a contradiction,

if a is even

Then $a=2\cdot q+1$. Note that
\[
\begin{aligned}
     2\cdot(6 \cdot k^2 + 6 \cdot k + 1) + 1&= a^2\\
     2\cdot(6 \cdot k^2 + 6 \cdot k + 1) + 1&= (2\cdot q+1)^2\\
     2\cdot(6 \cdot k^2 + 6 \cdot k + 1) + 1&= 4\cdot q^2+4\cdot q+1\\
     6 \cdot k^2 + 6 \cdot k + 1 &= 4\cdot q^2+4\cdot\\
     2\cdot (3 \cdot k^2 + 3 \cdot k) + 1 &= 2\cdot (2\cdot q^2+2)\cdot
\end{aligned}
\]

The left side is a odd number and the right side is even number, this is a contradiction.

Therefore $\sqrt{3}$ is irrational . $\blacksquare$
\end{document}